\documentclass[11pt,letterpaper]{article}
\usepackage[margin=1in]{geometry}
\usepackage[T1]{fontenc}
\usepackage{lmodern,amsmath,amssymb,amsthm,microtype}
\usepackage[hidelinks]{hyperref}
\hypersetup{pdftitle={A formal homogeneous Jacobian ideal separator},pdfsubject={Full operator separation under finite-dimensional gradient quotient}}
\newtheorem{theorem}{Theorem}[section]
\newtheorem{lemma}[theorem]{Lemma}
\newtheorem{proposition}[theorem]{Proposition}
\DeclareMathOperator{\Der}{Der}
\DeclareMathOperator{\End}{End}
\DeclareMathOperator{\Hess}{Hess}
\DeclareMathOperator{\adj}{adj}
\title{A formal homogeneous Jacobian ideal separator}
\author{}
\date{12 October 2026}
\begin{document}
\maketitle
\begin{abstract}
We formalize an explicit second-order differential operator outside the full
algebra generated by multiplication operators and all derivations for every
ordinary homogeneous polynomial of degree at least two in at least three
variables over a characteristic-zero field, assuming its gradient quotient is
finite-dimensional. An elementary determinant comparison with pure powers
proves the required Hessian nonmembership inside Lean, removing the classical
socle theorem previously left as an external mathematical input. The complete
sixteen-module proof and an expanded specification pass an independent
fresh-output rebuild. This extends the formal scope of the fixed-quartic
certificate; the operator and Jacobian-ideal stability retain their prior
attributions. Mathematical priority is not established. The weighted case and
the unrestricted Nakai conjecture are not resolved here.
\end{abstract}

\section{Exact scope and operator}
Let \(k\) be a field of characteristic zero, let
\(P=k[x_1,\ldots,x_n]\), and write
\[
 f_i=\partial_i f,\quad J=(f_1,\ldots,f_n),\quad
 H=(\partial_i\partial_j f)_{ij},\quad U=\adj(H),\quad h=\det H.
\]
For a \(k\)-algebra \(A\), put
\[
 \Delta_k(A)=k\langle m_a,\delta : a\in A,\ \delta\in\Der_k(A)\rangle
 \subseteq\End_k(A),\qquad m_a(b)=ab.
\]
There is no bound on the number of derivation factors in this algebra.
Differential order at most two means
\begin{equation}\label{eq:order}
 [[[T,m_a],m_b],m_c]=0\qquad(a,b,c\in A).
\end{equation}
This is the usual commutator characterization on actual endomorphisms.

\begin{theorem}\label{thm:main}
Suppose \(n\ge3\), \(d\ge2\), \(f\) is homogeneous of degree \(d\), and
\(\dim_k(P/J)<\infty\). Fix \(r\in\{1,\ldots,n\}\). Define
\begin{align}
 B_{ij}&=U_{ir}x_j+x_iU_{jr}-x_rU_{ij},\label{eq:B}\\
 Q_r(p)&=\sum_{i,j}B_{ij}\partial_i\partial_jp
 +(n-2)(d-1)\sum_iU_{ir}\partial_ip.\label{eq:Q}
\end{align}
Then \(Q_r\) descends to an actual \(k\)-linear endomorphism \(\overline Q_r\)
of \(A=P/(f)\), satisfies \eqref{eq:order}, and
\(\overline Q_r\notin\Delta_k(A)\).
\end{theorem}

The quotient finiteness hypothesis is expressed as \texttt{Module.Finite}
over the field. No Hessian nonmembership,
regular-sequence, or socle assertion is assumed. The algebraic operator theorem
is formalized for arbitrary finite index types. Geometric translations of the
finiteness hypothesis are not additional formal statements of this package.

\paragraph{Prior work.}
The operator \eqref{eq:Q} is twice the specified ordinary-homogeneous lift in
Yau--Zhu--Zuo \cite[eqs.~(3.1), (3.3), (5.1), (5.6), Cor.~5.5]{YZZ}, including
its first-order term. In three variables, the corresponding operator of
Diethorn et al.\ is \(Q_r/(d-1)^2\)
\cite[Sections 4.1.3 and 4.5, eqs.~(33), (35), Theorem~4.5.1]{DJM}; older
three-variable full-ring results also predate this work. Jacobian-ideal
stability is recorded in Li--Xiao--Zuo \cite[Theorem~2.7]{LXZ}.
Their Remark~1.1 explicitly distinguishes bounded-length generation from
intersection with the full generated algebra. The introduction and main theorem
of \cite{YZZ} formulate an order-by-order second-order obstruction; we do not
use that obstruction alone to infer unrestricted-word separation.
The purpose here is a complete checked proof of the stated algebraic theorem,
not a priority claim for its ingredients or assembled conclusion.

\section{An elementary Hessian nonmembership proof}
This section supplies the step that was previously imported from classical
complete-intersection socle theory. Determinantal nonmembership and
transformation methods are classical; see, for example,
Simon--Strooker \cite[Propositions 1.2 and 2.1]{SS} and
Eisenbud--Ulrich \cite[Theorem~7.1]{EU}. The proof below is chosen to fit a small,
explicit formal dependency graph.

\begin{lemma}\label{lem:wedge}
Let \(M\) be a square matrix over a commutative ring, \(x\) a column vector,
and \(I\) an ideal containing every entry of \(Mx\). For all \(r,p,q\),
\[
 x_q\adj(M)_{pr}-x_p\adj(M)_{qr}\in I.
\]
\end{lemma}
\begin{proof}
Replace row \(r\) of \(M\) by the basis row \(e_p^T\), obtaining \(L\).
Then \(\det L=\adj(M)_{pr}\), \((Lx)_r=x_p\), and
\(\adj(L)_{qr}=\adj(M)_{qr}\), because the latter cofactor deletes the
replaced row. The \(q\)-th entry of \(\adj(L)Lx=\det(L)x\) gives
\[
 x_q\adj(M)_{pr}-x_p\adj(M)_{qr}
 =\sum_{\ell\ne r}\adj(L)_{q\ell}(Mx)_\ell\in I.
\]
\end{proof}

\begin{lemma}\label{lem:comparison}
Suppose \(C,D\in M_n(P)\), every entry of both matrices is homogeneous of the
same degree \(m\), and \(Cx=Dx=y\). Then
\(\det C-\det D\in(y_1,\ldots,y_n)\).
\end{lemma}
\begin{proof}
If \(u=(u_j)\) is a row with entries homogeneous of degree \(m\) and
\(\sum_j u_jx_j=0\), Euler's identity and the derivative of this relation give,
for any coefficients \(c_j\),
\begin{equation}\label{eq:homotopy}
 (m+1)\sum_jc_ju_j
 =\sum_{i,j}(\partial_i u_j)(x_ic_j-x_jc_i).
\end{equation}
Indeed, the coefficient of \(c_j\) on the right is
\(\sum_i x_i\partial_i u_j-\sum_i x_i\partial_j u_i
 =mu_j-(-u_j)\).
For a matrix \(M\) with \(Mx=y\), take \(c_j=\adj(M)_{jr}\).
Every parenthesized term in \eqref{eq:homotopy} belongs to \((y)\) by
Lemma~\ref{lem:wedge}. Since \(m+1\) is invertible in \(k\), the row determinant
\(\det(M[r\leftarrow u])=\sum_j\adj(M)_{jr}u_j\) belongs to \((y)\).
Replace the rows of \(D\) successively by those of \(C\).
Each row difference is such a homogeneous syzygy, and every intermediate
matrix still has product \(y\) with \(x\). Row linearity and a finite sum of
these determinant differences prove the claim.
\end{proof}

\begin{proposition}\label{prop:detector}
Let \(g_1,\ldots,g_n\) be homogeneous of the same positive degree \(s\),
let \(J=(g_1,\ldots,g_n)\), and suppose \(P/J\) is finite-dimensional.
If \(g=Ex\), where the entries of \(E\) are homogeneous of degree \(s-1\),
then \(\det E\notin J\).
\end{proposition}
\begin{proof}
Each class \([x_i]\in P/J\) is integral over \(k\). Choose a monic polynomial
\(p_i(t)\) annihilating it. Then \(p_i(x_i)\in J\). The ideal \(J\) is
homogeneous, and the top homogeneous component of \(p_i(x_i)\) is
\(x_i^{\deg p_i}\); hence some power of each variable lies in \(J\).
Choose a common exponent \(N\ge s\), \(N>0\), at least as large as all those
exponents. Thus \(x_i^N=\sum_j b_{ij}g_j\).

Taking degree \(N\) components replaces every \(b_{ij}\) by its component of
degree \(N-s\), without changing this identity. The fact used here is
\[
 [bg]_N=[b]_{N-s}g\quad\text{if }g\text{ is homogeneous of degree }s
 \text{ and }N\ge s.
\]
Let \(B=(b_{ij})\) denote these homogeneous coefficients. Then
\(BEx=(x_i^N)_i\). The entries of \(BE\) and of
\(D=\operatorname{diag}(x_i^{N-1})\) have common degree \(N-1\).
Lemma~\ref{lem:comparison}, with \(I=(x_1^N,\ldots,x_n^N)\), gives
\begin{equation}\label{eq:transformation}
 \det B\det E-\prod_i x_i^{N-1}\in I.
\end{equation}
Also \(Bg=(x_i^N)_i\), so the adjugate identity gives
\(\det B\,g=\adj(B)(x_i^N)_i\). Consequently \(\det B\,J\subseteq I\).
If \(\det E\in J\), equation~\eqref{eq:transformation} would imply
\(\prod_i x_i^{N-1}\in I\). This is impossible: none of the generators
\(x_i^N\) divides that monomial. The monomial-ideal membership criterion
is checked for the actual polynomial ideal in the formal proof.
\end{proof}

Apply Proposition~\ref{prop:detector} to \(g_i=f_i\), \(s=d-1\), and
\(E=(d-1)^{-1}H\). Differentiated Euler gives \(Hx=(d-1)\nabla f\), so
\(Ex=\nabla f\). The proposition yields
\begin{equation}\label{eq:hessian}
 h=\det H\notin J.
\end{equation}
Thus the full Hessian nonmembership conclusion, rather than a classical socle
theorem used as an assumption, is now a kernel-checked consequence of finiteness
and homogeneity.

\section{Descent and unrestricted word separation}
The generic operator modules establish the following exact polynomial identities:
\begin{align}
 Q_r(fp)&=f\left(Q_r(p)+2d\sum_iU_{ir}\partial_i p\right)
       &&(p\in P),\label{eq:descent}\\
 Q_r(f_r)&=(nd-n-d)h-(d-1)\sum_i(\partial_rU_{ir})f_i.\label{eq:escape}
\end{align}
Equation~\eqref{eq:descent} proves descent on every representative.
The expression \eqref{eq:Q} is a sum of polynomial multiples of compositions
of at most two partial derivations; the formal commutator proof gives
\eqref{eq:order} on \(P\) and transports it to \(P/(f)\).
The exact first-order term is included throughout.

Euler's identity gives \(f\in J\). Every derivation \(\delta\) of \(A\) can be
lifted, by lifting its values on the variables, to
\(\widetilde\delta=\sum_i u_i\partial_i\) on \(P\), with
\(\widetilde\delta(f)=vf\). Differentiating gives
\[
 \widetilde\delta(f_j)
 =(\partial_jv)f+vf_j-\sum_i(\partial_ju_i)f_i\in J.
\]
Hence every actual quotient derivation preserves the image \(\overline J\)
of \(J\) in \(A\). Multiplications preserve it too, so the entire subalgebra
\(\Delta_k(A)\) preserves \(\overline J\). This includes arbitrary finite
words and all cancellations among them; no classification or selected
list of derivation generators is used.

Since
\[
 nd-n-d=(n-2)(d-1)+(d-2)>0\quad(n\ge3,\ d\ge2),
\]
this integer is nonzero in \(k\). Equations~\eqref{eq:hessian} and
\eqref{eq:escape} imply \(Q_r(f_r)\notin J\).
Because \((f)\subseteq J\), its quotient class is outside \(\overline J\).
Thus \(\overline Q_r\) fails to preserve an ideal preserved by every element
of \(\Delta_k(A)\), proving Theorem~\ref{thm:main}.

\section{Formal verification and limits}
The checked final theorem is
\begin{center}
\texttt{NakaiHomogeneousFiniteExpected.}\\
\texttt{expected\_finite\_gradient\_full\_operator\_separation}.
\end{center}
Its expanded statement writes the polynomial operator, literal gradient
quotient, actual hypersurface quotient, triple commutators, and unrestricted
\texttt{Algebra.adjoin} of all multiplication maps and all quotient derivations.
The only geometric-algebraic input is \texttt{Module.Finite} of the gradient
quotient. The determinant and finite-quotient lemmas prove the missing
nonmembership implication, without assuming a socle result.

All sixteen project modules were independently rebuilt into a fresh output
directory with Lean 4.32.2 and the following exact pins:
\begin{quote}\small
Lean: \texttt{f3b06c705e6c85f5314019d5d3baab0fec5b580c}\\
Mathlib: \texttt{905b95818eb32af7874a58b427f50c1711a5e96c}.
\end{quote}
One unchanged pinned upstream module, \texttt{MvPolynomial/EulerIdentity}, was compiled privately using
already cached dependencies. No previous project object files were used.
Every module exited successfully; the final axiom report lists only
\texttt{propext}, \texttt{Classical.choice}, and \texttt{Quot.sound}.
Exact source hashes, logs, guards, and the reproducible dependency order are
included with the source package.

This addendum supersedes the earlier fixed-quartic note's limitation that the
general ordinary-homogeneous statement was only a written theorem.
It does not alter that note's fifteen-source concrete certificate or its
operator attribution. It does not formalize a weighted-homogeneous extension,
settle the general nonhomogeneous case, or establish historical novelty of the
full-word argument. Those distinctions remain explicit.

\begin{thebibliography}{9}
\bibitem{YZZ}
S.~S.-T. Yau, Q. Zhu, and H. Zuo,
\emph{Nakai conjecture for isolated weighted homogeneous hypersurface singularities},
arXiv:2604.24508v4, 8 October 2026.
\url{https://arxiv.org/html/2604.24508v4}.
\bibitem{DJM}
R.~N. Diethorn, J. Jeffries, C. Miller, N. Packauskas, J. Pollitz,
H. Rahmati, and S. Vassiliadou,
\emph{Resolutions of differential operators of low order for an isolated hypersurface singularity},
arXiv:2209.13110v3.
\url{https://arxiv.org/html/2209.13110v3}.
\bibitem{LXZ}
R. Li, Z. Xiao, and H. Zuo,
\emph{The Nakai Conjecture for isolated hypersurface singularities of modality $\le2$},
arXiv:2502.04672v1.
\url{https://arxiv.org/pdf/2502.04672v1}.
\bibitem{SS}
A.-M. Simon and J.~R. Strooker,
\emph{Complete intersections of dimension zero: variations on a theme of Wiebe},
arXiv:math/0703880.
\url{https://arxiv.org/abs/math/0703880}.
\bibitem{EU}
D. Eisenbud and B. Ulrich,
\emph{Duality and socle generators for residual intersections},
Journal f\"ur die reine und angewandte Mathematik \textbf{756} (2019), 183--226.
\url{https://par.nsf.gov/servlets/purl/10282767}.
\end{thebibliography}
\end{document}
